% Introduction: the question, the approach, and the three contributions.
\section{Introduction}
\label{sec:intro}

Alignment theory has long asked what an advanced AI system should ultimately aim at. One of the best-known answers is
Coherent Extrapolated Volition \citep{yudkowsky2004cev}:
\begin{quote}\small
``our wish if we knew more, thought faster, were more the people we wished we were, had grown up farther together;
where the extrapolation converges rather than diverges, where our wishes cohere rather than interfere; extrapolated as
we wish that extrapolated, interpreted as we wish that interpreted.''
\end{quote}
CEV was written as the goal of a hypothetical superintelligence, not as instructions for a chat assistant. Today,
though, the cheapest way to give a deployed language model a value specification is to put it in the system prompt.
Practitioners do this with short role descriptions (``You are a helpful assistant''), long safety preambles
\citep{touvron2023llama2,askell2021general} and published constitutions \citep{bai2022constitutional,
anthropic2026constitution}. Whether a value text from alignment theory makes a model behave better, and on which
dimensions, is an empirical question. To our knowledge it has not been tested for CEV.

We ask a simple question. If a model is handed CEV verbatim as its system prompt, does it behave in a more aligned way
than with a classic assistant prompt, on evaluations that measure what CEV is meant to achieve? We place Section~3 of
\citet{yudkowsky2004cev} in the system prompt: the definition and the author's gloss of every clause, minimally
rewritten into the second person. We compare it with the classic Llama~2 assistant prompt \citep{touvron2023llama2}
and with no prompt, on four models from four developers. A length-matched condition adds an equally long, value-neutral
text to the classic prompt, which separates the effect of CEV's content from the effect of a longer prompt. The six
evaluations were chosen clause by clause:
\begin{itemize}
  \item human wellbeing and humaneness \citep{humanebench2025};
  \item human agency, including deferring decisions and not manipulating values \citep{sturgeon2025humanagencybench};
  \item manipulative dark patterns \citep{kran2025darkbench};
  \item corrigibility, power-seeking and survival instinct \citep{perez2022discovering};
  \item instrumental convergence \citep{he2025instrumentaleval};
  \item blackmail in an agentic scenario \citep{lynch2025agentic}.
\end{itemize}

\paragraph{Contributions.}
\begin{itemize}
  \item \textbf{A controlled test of CEV as a system prompt.} We run five prompt conditions, including a no-prompt
  baseline and a length-matched control, on four models (GPT, Claude, Qwen, Llama) and six CEV-mapped evaluations:
  \RNSamples{} item-paired samples analysed with paired tests and multiplicity correction.
  \item \textbf{A mixed result, located in the mechanism.} CEV does not beat the classic prompt. Adding it improves
  corrigibility- and power-related answers. It lowers deference to users, which is not significant after
  correction for the benchmark as a whole but concentrated in deciding for users. The losses concentrate in answers that
  explicitly reason ``as CEV'', and a judge-blinding test attributes them to behaviour rather than to the judge.
  \item \textbf{An open, auditable harness.} The code is built on Inspect \citep{inspect2024,inspectevals2024}. It
  verifies every API request against its condition, checks the source texts verbatim against the originals, and
  generates every number in this paper from the logs.
\end{itemize}

Section~\ref{sec:related} reviews related work and Section~\ref{sec:method} describes the design.
Section~\ref{sec:results} reports the results, and Section~\ref{sec:robustness} tests whether they survive
multiplicity, missing data, model removal, provider routing and the judge. We discuss the mechanism in
Section~\ref{sec:discussion} and the limitations in Section~\ref{sec:limitations}.
